%============================================================
%  app_quantum.tex
%  Appendix E: Operators on Hilbert Spaces — Reference Summary
%
%  Tensor Formats in Banach Spaces
%  Antonio Falcó, CEU Universities
%
%  This appendix is a reference guide to the operator-theoretic
%  notions used in Part V. The full exposition with proofs is in
%  Chapters 12 and 13. This appendix contains:
%    - A quick-reference index to Chapters 12-13
%    - Self-contained statements of the foundational theorems
%      (spectral theorem, SVD, trace class)
%    - The tensor product identifications connecting operators
%      with the framework of Chapters 1-4
%    - Historical notes
%
%  Standard references: ReedSimon1972, BratteliRobinson1987, Ryan2002
%============================================================

\chapter{Operators on Hilbert Spaces: Reference Summary}
\label{app:quantum}

This appendix is a reference guide to the operator-theoretic vocabulary
used in Part~\ref{part:quantum}. The full mathematical development ---
with complete proofs and motivation --- is in
Chapters~\ref{chap:quantum_states} and~\ref{chap:lindblad}.
This appendix is intended for two purposes: as a \emph{quick reference}
for readers who are familiar with operator theory and want to locate a
specific result, and as a \emph{self-contained introduction} for
readers without prior background in quantum mechanics who wish to
read Part~\ref{part:quantum}.

Throughout this appendix, $H$ denotes a separable complex Hilbert space
of \emph{arbitrary} (possibly infinite) dimension.
Standard references: \cite{ReedSimon1972}, \cite{BratteliRobinson1987},
and~\cite{Ryan2002} for the trace-class tensor-product perspective.

%------------------------------------------------------------
\section{Quick-reference index to Part~\ref{part:quantum}}
\label{app:quantum:index}
%------------------------------------------------------------

\begin{center}
\renewcommand{\arraystretch}{1.35}
\begin{tabular}{@{}p{7.5cm}l@{}}
\toprule
\textbf{Concept} & \textbf{Location} \\
\midrule
Inner product, Riesz isomorphism $J: H \to H^*$
    & \S\ref{app:quantum:hilbert} \\
Dual $H^*$ vs.\ conjugate space $\overline{H}$
    & Remark~\ref{rem:app:dual_vs_conjugate} \\
Dirac notation (not used; comparison only)
    & Remark~\ref{rem:app:dirac} \\
Adjoint $A^*$, self-adjoint, positive, unitary operators
    & \S\ref{app:quantum:bounded} \\
Orthogonal projections and closed subspaces
    & Prop.~\ref{prop:split_equiv} \\
Rank-one operators $u \otimes \varphi \in \cL(H)$
    & \S\ref{app:quantum:bounded},
      Remark~\ref{rem:app:rank_one_tensor} \\
Compact operators, spectral theorem
    & \S\ref{app:quantum:compact},
      Thm.~\ref{thm:app:spectral} \\
Singular values, SVD
    & \S\ref{app:quantum:compact},
      Thm.~\ref{thm:app:svd} \\
Hilbert--Schmidt operators $\cT{2}(H)$,
    $\langle A,B\rangle_{\mathrm{HS}} = \Tr(A^*B)$
    & \S\ref{app:quantum:hs},
      Def.~\ref{def:app:hs} \\
$\cT{2}(H) \cong H \otimes_H H^*$ isometrically
    & Thm.~\ref{thm:app:hs_tensor},
      Remark~\ref{rem:app:hs_tensor} \\
Trace-class operators $\cT{1}(H)$, trace norm $\|\cdot\|_1$
    & \S\ref{app:quantum:trace},
      Def.~\ref{def:app:trace_class} \\
$\cT{1}(H) \cong H \otimes_\pi H^*$, \condINJ\ and $\|\cdot\|_1$
    & Remark~\ref{rem:app:trace_projective} \\
Partial trace $\TrS{\alpha^c}(\rho)$
    & \S\ref{app:quantum:trace} \\
Density operators $\Dstate{H}$, pure/mixed states
    & \S\ref{app:quantum:states},
      Def.~\ref{def:app:density} \\
$\supp(\rho_\alpha) = \Umin{\alpha}{v}$
    — the fundamental connection
    & Thm.~\ref{thm:qs:minimal_partial_trace}
      in Ch.~\ref{chap:quantum_states} \\
Rank stratification $\Dstate{H} = \bigsqcup_r \Drank{r}{H}$
    & Prop.~\ref{prop:qs:rank_tucker}
      in Ch.~\ref{chap:quantum_states} \\
Manifold structure of $\Drank{r}{H}$, tangent space
    & Thm.~\ref{thm:qs:manifold}
      in Ch.~\ref{chap:quantum_states} \\
Bounded projector $\hat{P}_\rho: \cTsa{1}(H) \to T_\rho\Drank{r}{H}$
    & Prop.~\ref{prop:qs:projector}
      in Ch.~\ref{chap:quantum_states} \\
Lindblad (GKSL) equation
    & \S\ref{sec:lind:equation}
      in Ch.~\ref{chap:lindblad} \\
Well-posedness of projected Lindblad dynamics
    & Thm.~\ref{thm:lind:wellposed}
      in Ch.~\ref{chap:lindblad} \\
Kinematic equations in Stiefel-core coordinates
    & \S\ref{sec:lind:kinematic}
      in Ch.~\ref{chap:lindblad} \\
\bottomrule
\end{tabular}
\end{center}

%------------------------------------------------------------
\section{Complex Hilbert spaces and the dual}
\label{app:quantum:hilbert}
%------------------------------------------------------------

A \emph{complex Hilbert space} $H$ carries an inner product
$\langle \cdot, \cdot\rangle: H \times H \to \CC$, conjugate-linear
in the first argument and linear in the second, positive definite, and
complete. The Riesz representation theorem gives the isometric
conjugate-linear isomorphism
\begin{equation}
\label{eq:app:riesz}
    J: H \longrightarrow H^*,
    \qquad J(v) = \langle v, \cdot\rangle,
\end{equation}
where $H^* = \cL(H, \CC)$ is the topological dual. This appendix works
throughout with $H^*$ (defined by the norm alone), not the conjugate
space $\overline{H}$ (defined by the linear structure). The two are
isomorphic as Banach spaces via $J$, but not canonically as complex
Banach spaces.

\begin{remark}[Dual vs.\ conjugate space]
\label{rem:app:dual_vs_conjugate}
In finite dimensions $H^* \cong H \cong \overline{H}$ as complex vector
spaces. In infinite dimensions, the Riesz map $J: H \to H^*$ is
conjugate-linear (not linear), making $H$ and $H^*$ non-isomorphic
as complex Banach spaces. The identification $\cT{1}(H) \cong H
\otimes_\pi H^*$ (Remark~\ref{rem:app:trace_projective}) requires
$H^*$ as the second factor to preserve the $\CC$-bilinearity of the
tensor product; using $\overline{H}$ instead introduces a
conjugate-linear identification that is inconsistent with the
Banach-space tensor framework of Chapter~\ref{chap:algtensor}.
\end{remark}

\begin{remark}[Dirac notation]
\label{rem:app:dirac}
In the physics literature, $v \in H$ is written $|v\rangle$ (a
\emph{ket}) and $J(v) \in H^*$ is written $\langle v|$ (a \emph{bra}),
so $\langle u, v\rangle_H = \langle u | v\rangle$. The rank-one operator
$u \otimes J(v)$ is written $|u\rangle\langle v|$. This monograph uses
exclusively the standard mathematical notation: $\langle u, v\rangle$
for the inner product, $J(v) \in H^*$ for the Riesz dual, and
$u \otimes \varphi \in \cL(H)$ for rank-one operators.
\end{remark}

%------------------------------------------------------------
\section{Bounded operators, adjoint, and special classes}
\label{app:quantum:bounded}
%------------------------------------------------------------

$\cL(H)$ with the operator norm is a C*-algebra. The adjoint
$A^* \in \cL(H)$ of $A$ is characterised by
$\langle A^*u, v\rangle = \langle u, Av\rangle$ for all $u, v \in H$,
with $\|A^*A\| = \|A\|^2$. Key classes:
self-adjoint ($A = A^*$), positive ($A \geq 0$: $\langle u, Au\rangle
\geq 0$), strictly positive ($A > 0$), unitary ($A^*A = AA^* = I$),
orthogonal projection ($A = A^* = A^2$). Every closed subspace $S
\subset H$ has an orthogonal projection $\PS{S}$ and
$H = S \oplus S^\perp$ (cf.\ Proposition~\ref{prop:finitedim_split},
which holds for \emph{all} closed subspaces in Hilbert spaces, not
just finite-dimensional ones).

For $u \in H$ and $\varphi \in H^*$, the rank-one operator
$u \otimes \varphi \in \cL(H)$ is defined by
$(u \otimes \varphi)(w) \coloneqq \varphi(w)\, u$.

\begin{remark}[Rank-one operators and the tensor product $H \otimes H^*$]
\label{rem:app:rank_one_tensor}
The finite-rank operators form the algebraic tensor product
$H \otimes_a H^*$ inside $\cL(H)$. The map $(u, \varphi) \mapsto
u \otimes \varphi$ is $\CC$-bilinear, making this consistent with the
tensor product framework of Chapter~\ref{chap:algtensor}. The
Hilbert--Schmidt completion gives $\cT{2}(H) \cong H \otimes_H H^*$
and the trace-class completion gives $\cT{1}(H) \cong H \otimes_\pi H^*$;
see Sections~\ref{app:quantum:hs}--\ref{app:quantum:trace}.
\end{remark}

%------------------------------------------------------------
\section{Compact operators, spectral theorem, and SVD}
\label{app:quantum:compact}
%------------------------------------------------------------

$A \in \cL(H)$ is \emph{compact} if it maps the unit ball to a
relatively compact set. The set $\cK(H)$ of compact operators is a
closed two-sided ideal, equal to the norm-closure of finite-rank
operators (in separable $H$).

\begin{theorem}[Spectral theorem for compact self-adjoint operators,
{\cite[Thm.\ VI.16]{ReedSimon1972}}]
\label{thm:app:spectral}
Let $A \in \cK(H)$ with $A = A^*$. Then there exist an orthonormal
sequence $\{u_k\}$ in $H$ and real $\lambda_k \to 0$ with
\begin{equation}
\label{eq:app:spectral}
    A = \sum_{k} \lambda_k\, u_k \otimes J(u_k),
\end{equation}
convergent in operator norm. The $\lambda_k$ are the non-zero
eigenvalues of $A$.
\end{theorem}

The \emph{singular values} of $A \in \cK(H)$ are the eigenvalues
$s_k(A) \geq 0$ of $|A| = (A^*A)^{1/2}$.

\begin{theorem}[Singular value decomposition]
\label{thm:app:svd}
For $A \in \cK(H)$ with non-zero singular values $(s_k(A))$, there
exist orthonormal sequences $\{u_k\}$, $\{v_k\}$ in $H$ with
\begin{equation}
\label{eq:app:svd}
    A = \sum_k s_k(A)\, v_k \otimes J(u_k),
\end{equation}
convergent in operator norm. For $d=2$ and $v \in H_1 \otimes_H H_2$,
this is the Schmidt decomposition; the Tucker rank $r_1(v) = r_2(v)$
equals the number of non-zero singular values
(Remark~\ref{rem:qs:schmidt} in Chapter~\ref{chap:quantum_states}).
\end{theorem}

%------------------------------------------------------------
\section{Hilbert--Schmidt operators}
\label{app:quantum:hs}
%------------------------------------------------------------

\begin{definition}
\label{def:app:hs}
$A \in \cL(H)$ is \emph{Hilbert--Schmidt} if
$\|A\|_2^2 \coloneqq \sum_k \|Ae_k\|^2 = \Tr(A^*A) < \infty$
for some (equivalently every) orthonormal basis $\{e_k\}$.
The space $\cT{2}(H)$ with $\langle A, B\rangle_{\mathrm{HS}}
\coloneqq \Tr(A^*B)$ is a Hilbert space, a two-sided ideal in
$\cL(H)$, with $\|A\| \leq \|A\|_2 = (\sum_k s_k(A)^2)^{1/2}$.
\end{definition}

\begin{theorem}[$\cT{2}(H) \cong H \otimes_H H^*$]
\label{thm:app:hs_tensor}
There is a canonical isometric isomorphism of Hilbert spaces
\begin{equation}
\label{eq:app:hs_iso}
    \cT{2}(H) \cong H \otimes_H H^*,
\end{equation}
under which the elementary tensor $u \otimes \varphi \in H \otimes_a H^*$
corresponds to the rank-one operator $w \mapsto \varphi(w)\, u$.
\end{theorem}

\begin{remark}[Why $H^*$ and not $\overline{H}$]
\label{rem:app:hs_tensor}
The identification~\eqref{eq:app:hs_iso} uses $H^*$ to preserve
$\CC$-bilinearity. In finite dimensions (or when $H = \overline{H}$ is
identified via $J$) one also sees $\cT{2}(H) \cong H \otimes_H \overline{H}$
in the literature, but this conflates a linear and a conjugate-linear
isomorphism. The present convention is consistent with the projective
tensor norm identification in the next section.
\end{remark}

%------------------------------------------------------------
\section{Trace-class operators}
\label{app:quantum:trace}
%------------------------------------------------------------

\begin{definition}
\label{def:app:trace_class}
$A \in \cL(H)$ is \emph{trace class} if
$\|A\|_1 \coloneqq \Tr|A| = \sum_k s_k(A) < \infty$.
The space $\cT{1}(H)$ with $\|\cdot\|_1$ is a Banach space (not Hilbert),
a two-sided ideal in $\cL(H)$, with $\cT{1}(H) \subset \cT{2}(H) \subset
\cK(H)$ and $\|A\| \leq \|A\|_2 \leq \|A\|_1$.
The \emph{trace} $\Tr(A) \coloneqq \sum_k \langle e_k, Ae_k\rangle$
is absolutely convergent and basis-independent.
The duality $\cT{1}(H)^* \cong \cL(H)$ holds via $\langle A, X\rangle
= \Tr(A^*X)$.
\end{definition}

\begin{remark}[$\cT{1}(H) \cong H \otimes_\pi H^*$ and \condINJ]
\label{rem:app:trace_projective}
The identification
\begin{equation}
\label{eq:app:tc_iso}
    \cT{1}(H) \cong H \otimes_\pi H^*
\end{equation}
holds isometrically~\cite{Grothendieck1953,Ryan2002}. On $H \otimes_a H^*$:
\begin{equation*}
    \|\cdot\|_\varepsilon \leq \|\cdot\|_H \leq \|\cdot\|_1 = \|\cdot\|_\pi,
\end{equation*}
where $\|\cdot\|_H$ is the Hilbert--Schmidt norm and $\|\cdot\|_\varepsilon$
the injective norm. Both $\|\cdot\|_H$ and $\|\cdot\|_\pi$ satisfy
condition~\condINJ\ of Definition~\ref{def:crossnorm}. The trace norm
$\frac{1}{2}\|\rho - \sigma\|_1$ equals the total variation distance
between quantum states $\rho$ and $\sigma$ — the maximum probability
of distinguishing them by any measurement — which is why it is the
natural norm for state estimation and compression.
\end{remark}

The \emph{partial trace} $\TrS{\alpha^c}(A) \in \cT{1}(H_\alpha)$,
for $A \in \cT{1}(H_1 \otimes_H H_2)$, is characterised by
$\Tr_{H_1}(B\, \TrS{2}(A)) = \Tr_H((B \otimes I_2)\, A)$
for all $B \in \cL(H_1)$.

%------------------------------------------------------------
\section{Density operators and quantum states}
\label{app:quantum:states}
%------------------------------------------------------------

\begin{definition}
\label{def:app:density}
A \emph{density operator} is $\rho \in \cTsa{1}(H) \coloneqq
\{A \in \cT{1}(H) : A = A^*\}$ with $\rho \geq 0$ and $\Tr(\rho)=1$.
The \emph{state space} is $\Dstate{H} \coloneqq \{\rho \in \cTsa{1}(H)
: \rho \geq 0,\, \Tr(\rho)=1\}$.
\end{definition}

$\Dstate{H}$ is convex and trace-norm closed, but not globally a
manifold (rank singularities, boundary phenomena).
A state is \emph{pure} if $\mathrm{rank}(\rho)=1$ (equivalently
$\rho = u \otimes J(u)$ for a unit vector $u \in H$) and \emph{mixed}
otherwise.

\medskip\noindent
The \emph{rank-$r$ stratum} $\Drank{r}{H} \coloneqq \{\rho \in \Dstate{H}:
\mathrm{rank}(\rho)=r\}$ is a smooth Banach manifold (Theorem~\ref{thm:qs:manifold}
in Chapter~\ref{chap:quantum_states}).
The stratification $\Dstate{H} = \bigsqcup_r \Drank{r}{H} \cup
\mathcal{D}_\infty(H)$ mirrors the Tucker rank decomposition
$\VD = \bigsqcup_\fr \Tucker{\fr}{\VD}$ of
Theorem~\ref{thm:tucker_decomposition_full}.

\medskip\noindent
The fundamental connection with minimal subspaces
(Theorem~\ref{thm:qs:minimal_partial_trace}) is:
for $v \in \bigotimes_{j=1}^d H_j$ with $\|v\|=1$,
\begin{equation}
\label{eq:app:min_supp}
    \supp\bigl(\TrS{\alpha^c}(v \otimes J(v))\bigr) = \Umin{\alpha}{v},
    \qquad
    \mathrm{rank}\bigl(\TrS{\alpha^c}(v \otimes J(v))\bigr) = r_\alpha(v).
\end{equation}
This is the cornerstone of Chapter~\ref{chap:quantum_states}.

%------------------------------------------------------------
\section{Historical notes}
\label{app:quantum:history}
%------------------------------------------------------------

\paragraph{Operator-theoretic foundations of quantum mechanics.}
Von Neumann~\cite{vonNeumann1932} established the functional-analytic
framework using density operators on Hilbert spaces. Schatten~\cite{Schatten1950}
introduced the Schatten $p$-classes $\cT{p}(H)$ with norm
$\|A\|_p = (\sum_k s_k(A)^p)^{1/p}$; the cases $p=1$ (trace class)
and $p=2$ (Hilbert--Schmidt) are central to Part~\ref{part:quantum}.

\paragraph{Tensor norms and the trace.}
The identification $\cT{1}(H) \cong H \otimes_\pi H^*$
(Remark~\ref{rem:app:trace_projective}) is due to
Grothendieck~\cite{Grothendieck1953}, connecting the trace norm to the
projective tensor norm of Chapter~\ref{chap:algtensor}. The
Hilbert--Schmidt identification $\cT{2}(H) \cong H \otimes_H H^*$
(Theorem~\ref{thm:app:hs_tensor}) is the non-commutative analogue of
the Fubini isomorphism $L^2(\Omega_1) \otimes_H L^2(\Omega_2) \cong
L^2(\Omega_1 \times \Omega_2)$ (cf.~\cite[Chapter~II]{ReedSimon1972}).

\paragraph{Minimal subspaces and partial traces.}
Equation~\eqref{eq:app:min_supp} makes precise a connection implicit in
the quantum information literature since Schmidt~(1907). Its systematic
treatment in the tensor-Banach-space framework is the contribution of
Falcó--Hackbusch~\cite{FH2012} and Acevedo--Falcó~\cite{AcevedoFalco2026}.
